% Generated draft: requires mathematical and source review.
\documentclass[11pt]{article}
\usepackage[margin=25mm]{geometry}
\usepackage{fontspec}
\setmainfont{Latin Modern Roman}
\usepackage{amsmath,amssymb,mathtools,unicode-math}
\setmathfont{Latin Modern Math}
\usepackage{enumitem,needspace,xurl,hyperref}
\setlength{\parindent}{0pt}
\setlength{\parskip}{5pt}
\setlength{\emergencystretch}{3em}
\newcommand{\N}{\mathbb{N}}
\newcommand{\Z}{\mathbb{Z}}
\newcommand{\Q}{\mathbb{Q}}
\newcommand{\R}{\mathbb{R}}
\newcommand{\C}{\mathbb{C}}
\newcommand{\F}{\mathbb{F}}
\newcommand{\A}{\mathbb{A}}
\newcommand{\PP}{\mathbb{P}}
\newcommand{\E}{\mathbb{E}}
\newcommand{\eps}{\varepsilon}
\newcommand{\dd}{\,\mathrm{d}}
\begin{document}
\section*{UnsolvedMath ID: 701749}
\textbf{Source problem code:} catalogue entry
\subsection{Definitions and mathematical statement}
Let \(L_{\mathrm{grp}}=\{\cdot,{}^{-1},1\}\) be the first-order language of groups. A group \(G\) is \emph{stable} if, for every \(L_{\mathrm{grp}}\)-formula \(\phi(\bar x,\bar y)\), there exists an integer \(n\geq 0\) such that, for every integer \(m\geq 1\) and tuples
\[
 a_i\in G^{|\bar x|},\qquad b_i\in G^{|\bar y|}\qquad(0\leq i<m),
\]
if
\[
 G\models\phi(a_i,b_j)\quad\Longleftrightarrow\quad i<j
 \qquad\text{for all }0\leq i,j<m,
\]
then \(m\leq n\). Thus \(G\) is \emph{non-stable} precisely when there is one such formula \(\phi\) for which these configurations exist for every integer \(m\geq 1\).

A group is \emph{hyperbolic} if it has a finite generating set whose Cayley graph, with every edge of length one and its path metric, has \(\delta\)-thin geodesic triangles for some real \(\delta\geq 0\): every side of every geodesic triangle lies in the closed \(\delta\)-neighborhood of the other two sides. A group is \emph{virtually free} if it contains a free subgroup of finite index. The supplied record does not specify whether finite generation is additionally intended for virtually free groups in part (b); the definition just given imposes no rank restriction on the free subgroup, and this convention requires review against the original source. Write
\[
 \mathrm{SL}_2(\Z)=\{A\in M_2(\Z):\det A=1\},
\]
with matrix multiplication.

The supplied background attributes to Sela the theorem that every torsion-free hyperbolic group is stable; torsion-free means that the identity is the only element of finite order. This attribution is supplied context and has not been independently verified here.

For each of the following, determine whether the displayed assertion holds:
\begin{enumerate}
\item[(a)] There exists a hyperbolic group that is non-stable.
\item[(b)] There exists a virtually free group that is non-stable, subject to the unresolved finite-generation convention above.
\item[(c)] The group \(\mathrm{SL}_2(\Z)\) is non-stable.
\end{enumerate}
\subsection{Short English statement}
Stability requires a finite bound, depending on the formula, on the lengths of configurations in which the formula holds exactly when the first index is smaller than the second. The problem asks whether non-stable groups occur among hyperbolic groups and virtually free groups, and whether \(\mathrm{SL}_2(\Z)\) is non-stable. The finite-generation convention in the virtually free question remains unspecified.
\subsection{Sources}
\begin{itemize}
\item Supplied local UnsolvedMath registry record, permanent ID \(701749\), exact source code \texttt{catalogue entry}, titled \emph{Kourovka Notebook Problem 20.91}. The record labels the problem \emph{open}; this is a supplied status, not an independently verified literature assessment.
\item \emph{The Kourovka Notebook (Unsolved Problems in Group Theory)}, original problem 20.91, proposed by R. Sklinos in 2022. The supplied background identifies the 20th issue (2022) and the 21st issue edition (2026), with identifier \texttt{arXiv:1401.0300v48}. The supplied collection description names E. I. Khukhro and V. D. Mazurov as editors of the Issues 1--21 collection (1965--2026).
\item The supplied statement attributes stability of all torsion-free hyperbolic groups to Sela. No separate publication citation for that theorem is supplied. The cited Notebook edition and attribution have not been independently checked here.
\end{itemize}
\end{document}
